\documentclass{article}
\usepackage{graphicx} % Required for inserting images

\title{ELEC8109 Project Proposal}
\author{Zongqi He, Qingyang Bao, Shengrong Zhang}
\date{\today}

\begin{document}

\maketitle

\section{Background}
3D Gaussian Splatting (3DGS)~\cite{3dgs} enables high-quality, real-time novel-view synthesis, providing an efficient scene representation for interactive virtual photography. However, its original formulation adopts an idealized pinhole camera model without explicitly accounting for finite-aperture defocus. To enable richer photographic controls, Cinematic Gaussians~\cite{cinematicgaussians}, DOF-GS~\cite{dofgs}, and DoF-Gaussian~\cite{dofgaussian} incorporate thin-lens imaging and differentiable defocus models into Gaussian rendering, supporting post-capture refocusing, aperture adjustment, and depth-of-field control.

Recent research further expands aperture modeling capabilities. Closed-Form Convolution~\cite{closedformconvolution} introduces efficient analytical defocus computation with support for polygonal apertures. EDoF-NeRF~\cite{edofnerf} incorporates coded apertures into radiance-field reconstruction, demonstrating the potential of spatially patterned pupil transmission to extend depth of field and improve reconstruction quality. Meanwhile, fNeRF~\cite{fnerf} highlights that rays passing through different aperture locations can encounter different visibility conditions, causing image-space convolution approximations to produce errors around partial occlusions and depth discontinuities.

Building on these advances, this project investigates programmable-aperture depth-of-field rendering for 3DGS, integrating aperture geometry, nonuniform pupil transmission, and depth-dependent defocus into a unified image formation model. We will analyze the errors of different defocus approximations around depth discontinuities and partial occlusions, and explore rendering methods that balance imaging accuracy with computational efficiency. From a computational photography perspective, our intended contribution is to combine aperture transmission control with an interactive 3D scene representation, enabling view-consistent post-capture aperture editing and refocusing while providing an efficient approach to controllable bokeh synthesis.
\section{Proposed Approach}
Our objective is to develop a programmable aperture model for 3DGS that enables post-capture rendering across arbitrary pupil configurations. We explicitly model spatial light transmission across the aperture and parameterize its coupling with depth-dependent defocus blur, extending conventional aperture size and focus distance controls to complex aperture geometries and non-uniform transmission patterns.

Building upon existing depth-of-field formulations, this research addresses how these richer optical controls can be tractably integrated into Gaussian rendering while maintaining strict consistency across viewpoints and scene depths. Specifically, we investigate the trade-off between physical accuracy and real-time efficiency under complex aperture patterns, with particular emphasis on resolving blur-visibility interactions near depth discontinuities.
\section{Experimental Setup \& Evaluation}
To validate our programmable aperture formulation, we evaluate rendering fidelity and post-capture controllability across both synthetic and real-world benchmarks:

\begin{itemize}
    \item \textbf{Datasets:} Following DoF-Gaussian~\cite{dofgaussian} and related depth-of-field literature, we utilize physically based synthetic scenes rendered via Blender with thin-lens optics under systematically varied focal planes, aperture patterns, and scene depths, providing all-in-focus pinhole ground truths. For real-world validation, we evaluate on standard defocused benchmarks (such as the Real-DoF dataset~\cite{realdof}), supplemented by custom multi-view captures captured under diverse aperture geometries (e.g., polygonal blades, coded masks) and pronounced depth discontinuities.
    \item \textbf{Baselines \& Metrics:} We benchmark against standard 3DGS~\cite{3dgs}, Cinematic Gaussians~\cite{cinematicgaussians}, DoF-Gaussian~\cite{dofgaussian}, and DOF-GS~\cite{dofgs}. Quantitative reconstruction quality is measured via PSNR, SSIM, and LPIPS against reference images. To verify runtime practicality, we evaluate training convergence time, peak GPU memory footprint, and interactive rendering throughput (FPS).
\end{itemize}
\section{Expected Contributions}
This project aims to deliver the following key contributions:
\begin{itemize}
    \item \textbf{Programmable Aperture Formulation:} A unified optical formulation that parameterizes spatially varying pupil transmission and arbitrary aperture geometries within the 3DGS pipeline, expanding post-capture photographic expressiveness beyond conventional circular thin-lens approximations.
    \item \textbf{Systematic Characterization \& Benchmark:} A comprehensive evaluation benchmark analyzing where analytical and ray-based approximations diverge from wave-optics ground truth, establishing concrete trade-offs between physical accuracy and rendering throughput.
\end{itemize}
Ultimately, this work bridges the gap between explicit neural representations and physical optical simulation, providing an efficient foundation for computational photography and cinematic post-processing.
\bibliographystyle{unsrt} %apacite
\bibliography{reference}
\end{document}
